\chapter{Conditional Expectation and Inequalities}
\label{ch:condexp}

Conditioning is the central operation of probability: it is how information
changes our estimates. This chapter extends conditional probability to random
variables, proves the tower property and the law of total variance, shows that
conditional expectation is the best predictor in mean square, and develops the
standard inequalities (Markov, Chebyshev, Jensen, Chernoff, Hoeffding) that turn
moments and transforms into probability bounds. It ends with symmetry arguments,
which solve many problems with no computation at all.

\section{Conditional distributions and expectation: the discrete case}

Let $X,Y$ be discrete. For $y$ with $p_Y(y)>0$ the \emph{conditional PMF} of $X$
given $Y=y$ is
\[
p_{X\given Y}(x\given y)=\Prob(X=x\given Y=y)=\frac{p_{X,Y}(x,y)}{p_Y(y)} .
\]
For fixed $y$ this is a PMF in $x$ (it is $\Prob(\cdot\given Y=y)$ restricted to
$\{X=x\}$, a probability measure by Theorem~\ref{thm:condmeasure}).

\begin{definition}[Conditional expectation]
For $y$ with $p_Y(y)>0$ define
$\psi(y)=\Exp[X\given Y=y]=\sum_xx\,p_{X\given Y}(x\given y)$, provided the series
converges absolutely. The random variable $\Exp[X\given Y]:=\psi(Y)$ is the
\emph{conditional expectation of $X$ given $Y$}.
\end{definition}

The notation hides an important point: $\Exp[X\given Y=y]$ is a \emph{number} (a
function of $y$), while $\Exp[X\given Y]$ is a \emph{random variable} (a function
of $Y$), whose value changes as $Y$ changes. It is ``the average of $X$ given what
$Y$ tells us''.

\begin{theorem}[Defining property and tower property]\label{thm:tower}
Let $X$ be integrable. Then for every bounded function $g$,
\[
\Exp\bigl[g(Y)\,X\bigr]=\Exp\bigl[g(Y)\,\Exp[X\given Y]\bigr].
\tag{$\ast$}
\]
In particular (take $g\equiv1$), the \emph{tower property}
$\Exp\bigl[\Exp[X\given Y]\bigr]=\Exp[X]$.
\end{theorem}
\begin{proof}
\[
\Exp[g(Y)X]=\sum_{x,y}g(y)\,x\,p_{X,Y}(x,y)=\sum_yg(y)p_Y(y)\sum_xx\frac{p_{X,Y}(x,y)}{p_Y(y)}
=\sum_yg(y)\psi(y)p_Y(y)=\Exp[g(Y)\psi(Y)].
\]
All rearrangements are justified by absolute convergence, because $g$ is bounded
and $X$ is integrable.
\end{proof}

\begin{theorem}[Properties of conditional expectation]\label{thm:condprops}
Assume all expectations exist.
\begin{enumerate}[label=(\roman*)]
\item (Linearity) $\Exp[aX+bZ\given Y]=a\Exp[X\given Y]+b\Exp[Z\given Y]$.
\item (Monotonicity) If $X\le Z$ then $\Exp[X\given Y]\le\Exp[Z\given Y]$.
\item (Taking out what is known) $\Exp[h(Y)X\given Y]=h(Y)\Exp[X\given Y]$.
\item (Independence) If $X$ and $Y$ are independent, $\Exp[X\given Y]=\Exp[X]$.
\item (Total law) $\Exp[X]=\sum_y\Exp[X\given Y=y]\,\Prob(Y=y)$.
\end{enumerate}
\end{theorem}
\begin{proof}
Each statement holds for each fixed value $Y=y$ because $\Exp[\,\cdot\given Y=y]$ is an
ordinary expectation under the probability measure $\Prob(\cdot\given Y=y)$, so
(i) and (ii) are linearity and monotonicity of expectation. (iii) Given $Y=y$ the
factor $h(Y)=h(y)$ is a constant. (iv) $p_{X\given Y}(x\given y)=p_X(x)$ by
independence. (v) is the tower property written out.
\end{proof}

\begin{nuance}{the total law is the law of total probability in disguise}
With $X=\indic_A$, we have $\Exp[X\given Y=y]=\Prob(A\given Y=y)$, and (v) becomes
$\Prob(A)=\sum_y\Prob(A\given Y=y)\Prob(Y=y)$: Theorem~\ref{thm:lotp}.
\end{nuance}

\subsection*{Best prediction in mean square}

\begin{theorem}[Conditional expectation is the best predictor]\label{thm:bestpred}
Let $\Exp[X^2]<\infty$. For every function $h$ with $\Exp[h(Y)^2]<\infty$,
\[
\Exp\bigl[(X-h(Y))^2\bigr]=\Exp\bigl[(X-\psi(Y))^2\bigr]+\Exp\bigl[(\psi(Y)-h(Y))^2\bigr]
\ \ge\ \Exp\bigl[(X-\psi(Y))^2\bigr],
\]
with equality iff $h(Y)=\psi(Y)$ almost surely. So $\psi(Y)=\Exp[X\given Y]$ is the
function of $Y$ minimising mean squared error.
\end{theorem}
\begin{proof}
Write $X-h(Y)=(X-\psi(Y))+(\psi(Y)-h(Y))$ and expand the square. The cross term is
$2\Exp[(X-\psi(Y))\,g(Y)]$ with $g=\psi-h$; by the defining property $(\ast)$,
$\Exp[g(Y)X]=\Exp[g(Y)\psi(Y)]$, so the cross term vanishes. (Boundedness of $g$
in $(\ast)$ is removed by a limiting argument using $\Exp[g(Y)^2]<\infty$;
see Chapter~\ref{ch:measure}.)
\end{proof}

Geometrically, $X-\Exp[X\given Y]$ is ``orthogonal'' to every function of $Y$, and
$\Exp[X\given Y]$ is the orthogonal projection of $X$ onto the space of functions
of $Y$. This viewpoint is the key to the abstract definition in
Chapter~\ref{ch:measure}.

\subsection*{The law of total variance}

Define the \emph{conditional variance}
$\Var(X\given Y)=\Exp\bigl[(X-\psi(Y))^2\given Y\bigr]$.

\begin{theorem}[Law of total variance]\label{thm:totalvar}
If $\Exp[X^2]<\infty$,
\[
\Var(X)=\Exp\bigl[\Var(X\given Y)\bigr]+\Var\bigl(\Exp[X\given Y]\bigr).
\]
\end{theorem}
\begin{proof}
Write $X-\mu=(X-\psi(Y))+(\psi(Y)-\mu)$ where $\mu=\Exp[X]=\Exp[\psi(Y)]$. Square
and take expectations. The cross term $2\Exp[(X-\psi(Y))(\psi(Y)-\mu)]$ is zero by
$(\ast)$ with $g=\psi-\mu$. The first square has expectation
$\Exp[\Exp[(X-\psi(Y))^2\given Y]]=\Exp[\Var(X\given Y)]$ by the tower property, and
the second is $\Var(\psi(Y))$.
\end{proof}

The two terms are the \emph{unexplained} variance (average spread of $X$ around its
prediction) and the \emph{explained} variance (spread of the prediction itself).

\begin{example}[Random sums, again]
Let $S=\sum_{i=1}^NX_i$ with $X_i$ i.i.d.\ (mean $\mu$, variance $\sigma^2$), $N$
independent of the $X_i$. Given $N=n$, $\Exp[S\given N]=N\mu$ and
$\Var(S\given N)=N\sigma^2$. The tower property and the law of total variance give
\[
\Exp[S]=\mu\Exp[N],\qquad\Var(S)=\sigma^2\Exp[N]+\mu^2\Var(N),
\]
exactly as in Theorem~\ref{thm:randomsum}, with no generating functions.
\end{example}

\section{The continuous case}

If $(X,Y)$ has joint density $f$, the \emph{conditional density} of $X$ given $Y=y$ is
\[
f_{X\given Y}(x\given y)=\frac{f(x,y)}{f_Y(y)},\qquad f_Y(y)>0,
\]
and $\Exp[X\given Y=y]=\int xf_{X\given Y}(x\given y)\,dx$. The defining
property $(\ast)$, the tower property, the total variance formula and
Theorems~\ref{thm:condprops} and \ref{thm:bestpred} hold with the same proofs, with
sums replaced by integrals (the interchange of integrals is justified by Fubini's
theorem, Chapter~\ref{ch:measure}). Example: for the bivariate normal,
Theorem~\ref{thm:bvn}(ii) gives $\Exp[Y\given X]=\rho X$ and $\Var(Y\given X)=1-\rho^2$,
and the total variance formula recovers $\Var(Y)=(1-\rho^2)+\rho^2=1$.

\begin{nuance}{conditioning on a probability-zero event is a definition, not a theorem}
For continuous $Y$ we condition on $\{Y=y\}$, which has probability zero, so
$\Prob(A\given Y=y)=\Prob(A\cap\{Y=y\})/\Prob(Y=y)$ is $0/0$. What we have
defined is a \emph{function} $\psi(y)$ that satisfies $(\ast)$; it is determined only
up to a set of $y$ of probability zero, and its value at any single point is
arbitrary. It also depends on \emph{how} the null event is approached. Take a
point uniformly distributed on a sphere and condition on ``the point lies on a given
great circle''. If the circle is regarded as the equator (condition on latitude
$=0$), the position along it is uniformly distributed. If the same kind of circle is
regarded as a meridian (condition on longitude fixed), the conditional density of the
latitude along it is proportional to the cosine of the latitude, which is \emph{not}
uniform. The same geometric set gets different conditional distributions depending on
the variable being conditioned on. This is the Borel--Kolmogorov paradox. The remedy is always to specify
the conditioning \emph{variable}, not just the event.
\end{nuance}

\section{First-step analysis}

The law of total expectation turns hard problems into linear equations: condition
on the outcome of the first step, and use the fact that the process restarts.

\begin{example}[Expected waiting time for HH and for HT]\label{ex:hhht}
A fair coin is tossed until a given pattern appears. Let $E$ be the expected number
of tosses. For \textbf{HH}: define states ``start'' ($E_0$) and ``last toss was H''
($E_1$). From the start, one toss is made; heads moves to state $1$, tails stays at
the start: $E_0=1+\tfrac12E_1+\tfrac12E_0$. From state $1$: heads ends the game,
tails returns to start: $E_1=1+\tfrac12\cdot0+\tfrac12E_0$. Substituting,
$E_0=1+\tfrac12(1+\tfrac12E_0)+\tfrac12E_0=\tfrac32+\tfrac34E_0$, so $E_0=6$.

For \textbf{HT}: $E_0=1+\tfrac12E_1+\tfrac12E_0$ as before, but from state $1$,
heads keeps us in state $1$ and tails ends: $E_1=1+\tfrac12E_1$, so $E_1=2$ and
$E_0=1+1+\tfrac12E_0$ gives $E_0=4$. The two patterns are equally likely in any
fixed window of tosses, yet HH takes longer on average: HH can overlap itself, so
its occurrences come in clumps, and the clumping increases the wait for the
\emph{first} occurrence.
\end{example}

\begin{example}[Waiting for $n$ heads in a row]\label{ex:nheads}
Let $E_n$ be the expected number of tosses to get $n$ consecutive heads. To get $n$
in a row we first need $n-1$ in a row (expected $E_{n-1}$ tosses), then one more
toss: heads (probability $\tfrac12$) completes the run; tails (probability
$\tfrac12$) sends us back to the start. Hence
\[
E_n=E_{n-1}+1+\tfrac12\cdot0+\tfrac12E_n\ \Longrightarrow\ E_n=2E_{n-1}+2,
\qquad E_0=0,
\]
and by induction $E_n=2^{n+1}-2$. For $n=2$ we get $6$, matching the previous example.
\end{example}

\begin{interview}{first-step analysis}
Whenever a process regenerates (coin sequences, random walks, dice-until-condition),
set up one equation per state: $E_s=1+\sum_{s'}\Prob(s\to s')E_{s'}$, with $E=0$ at
absorbing targets, and solve the linear system. We return to this systematically for
Markov chains in Chapter~\ref{ch:markov}.
\end{interview}

\section{Basic inequalities}

\begin{theorem}[Markov's inequality]\label{thm:markov}
If $X\ge0$ and $a>0$, then $\Prob(X\ge a)\le\Exp[X]/a$.
\end{theorem}
\begin{proof}
Pointwise, $a\indic_{\{X\ge a\}}\le X$ (if $X\ge a$ the left is $a\le X$; otherwise it
is $0\le X$). Take expectations: $a\Prob(X\ge a)\le\Exp[X]$.
\end{proof}

\begin{theorem}[Chebyshev's inequality]\label{thm:chebyshev}
If $\Exp[X^2]<\infty$ and $t>0$, then $\Prob(\abs{X-\mu}\ge t)\le\sigma^2/t^2$.
\end{theorem}
\begin{proof}
Apply Markov to the non-negative variable $(X-\mu)^2$ and the level $t^2$.
\end{proof}

Chebyshev says that at least $1-1/k^2$ of the mass lies within $k$ standard
deviations of the mean, for \emph{any} distribution with finite variance. It is
crude (for the normal the actual probability beyond $2\sigma$ is $0.046$, versus the
bound $0.25$) but universal.

\begin{theorem}[One-sided Chebyshev (Cantelli)]\label{thm:cantelli}
For $t>0$, $\Prob(X-\mu\ge t)\le\dfrac{\sigma^2}{\sigma^2+t^2}$.
\end{theorem}
\begin{proof}
For any $u\ge0$, $\{X-\mu\ge t\}\subseteq\{(X-\mu+u)^2\ge(t+u)^2\}$, so by Markov
\[
\Prob(X-\mu\ge t)\le\frac{\Exp[(X-\mu+u)^2]}{(t+u)^2}=\frac{\sigma^2+u^2}{(t+u)^2}.
\]
Choose $u=\sigma^2/t$ (which minimises the right side): the bound becomes
$\dfrac{\sigma^2+\sigma^4/t^2}{(t+\sigma^2/t)^2}=\dfrac{\sigma^2}{\sigma^2+t^2}$.
\end{proof}

\begin{theorem}[Jensen's inequality]\label{thm:jensen}
If $\varphi:\R\to\R$ is convex and $X$, $\varphi(X)$ are integrable, then
$\varphi(\Exp[X])\le\Exp[\varphi(X)]$. If $\varphi$ is strictly convex, equality holds
iff $X$ is almost surely constant.
\end{theorem}
\begin{proof}
A convex function has a supporting line at every point: with $\mu=\Exp[X]$ there is
a constant $c$ (a subgradient at $\mu$) such that
$\varphi(x)\ge\varphi(\mu)+c(x-\mu)$ for all $x$. Substitute $x=X$ and take
expectations: $\Exp[\varphi(X)]\ge\varphi(\mu)+c\cdot0=\varphi(\mu)$. For strictly
convex $\varphi$ the supporting line touches the graph only at $\mu$, so equality
forces $X=\mu$ almost surely.
\end{proof}

\begin{example}[Consequences of Jensen]
\begin{itemize}
\item $\varphi(x)=x^2$: $(\Exp[X])^2\le\Exp[X^2]$, i.e.\ $\Var(X)\ge0$.
\item $\varphi(x)=\abs x^p$, $p\ge1$: $\bigl(\Exp\abs X\bigr)^p\le\Exp\abs X^p$, so
$\norm{X}_r=(\Exp\abs X^r)^{1/r}$ is non-decreasing in $r\ge1$.
\item $\varphi(x)=-\ln x$ for $X>0$: $\Exp[\ln X]\le\ln\Exp[X]$, the AM--GM inequality
for expectations.
\item $\varphi(x)=e^{tx}$: $M_X(t)\ge e^{t\Exp[X]}$.
\end{itemize}
\end{example}

\begin{pitfall}{Jensen reverses for concave functions, and $\Exp[1/X]\ne1/\Exp[X]$}
For concave $\varphi$ the inequality reverses. For instance
$\Exp[1/X]\ge1/\Exp[X]$ for $X>0$ (since $1/x$ is convex). A common error in
finance and queueing is to replace $\Exp[g(X)]$ by $g(\Exp[X])$; Jensen shows the
sign of the error.
\end{pitfall}

\section{Exponential bounds: Chernoff and Hoeffding}

Chebyshev's polynomial tail $1/t^2$ is too weak for sums of many independent
variables. Applying Markov to an exponential gives far sharper bounds.

\begin{theorem}[Chernoff bound]\label{thm:chernoff}
For every $a$ and every $t>0$ in the domain of the MGF,
\[
\Prob(X\ge a)\le e^{-ta}M_X(t),\qquad\text{so}\qquad
\Prob(X\ge a)\le\inf_{t>0}e^{-ta}M_X(t).
\]
\end{theorem}
\begin{proof}
$\{X\ge a\}=\{e^{tX}\ge e^{ta}\}$ since $e^{tx}$ is increasing; apply Markov to the
non-negative variable $e^{tX}$.
\end{proof}

\begin{example}[Normal tail]
For $Z\sim N(0,1)$, $e^{-ta}M(t)=e^{-ta+t^2/2}$ is minimised at $t=a$, giving
$\Prob(Z\ge a)\le e^{-a^2/2}$ for $a>0$.
\end{example}

\begin{example}[Binomial tail via relative entropy]\label{ex:chernoffbin}
Let $X\sim\mathrm{Bin}(n,p)$ and $\alpha\in(p,1)$. Then $e^{-t\alpha n}M_X(t)=
\exp\bigl(n[-t\alpha+\ln(q+pe^t)]\bigr)$. Setting the derivative of the bracket to
zero, $-\alpha+pe^t/(q+pe^t)=0$, gives $e^t=\alpha q/((1-\alpha)p)$, and then
$q+pe^t=q/(1-\alpha)$. Substituting,
\[
-t\alpha+\ln(q+pe^t)=-\alpha\ln\frac\alpha p-(1-\alpha)\ln\frac{1-\alpha}q=-D(\alpha\|p),
\]
where $D(\alpha\|p)=\alpha\ln\frac\alpha p+(1-\alpha)\ln\frac{1-\alpha}{1-p}\ge0$ is the
\emph{relative entropy} (Kullback--Leibler divergence) between Bernoulli laws. Hence
\[
\Prob(X\ge\alpha n)\le e^{-nD(\alpha\|p)}.
\]
The tail probability decays \emph{exponentially} in $n$, with rate given by $D$.
(Cram\'er's theorem says this rate is exact to first order.)
\end{example}

\begin{lemma}[Hoeffding's lemma]\label{lem:hoeffding}
If $\Exp[X]=0$ and $a\le X\le b$ almost surely, then for every real $s$,
$\Exp[e^{sX}]\le\exp\bigl(s^2(b-a)^2/8\bigr)$.
\end{lemma}
\begin{proof}
Note $a\le0\le b$. By convexity of $x\mapsto e^{sx}$, for $x\in[a,b]$,
$e^{sx}\le\frac{b-x}{b-a}e^{sa}+\frac{x-a}{b-a}e^{sb}$. Taking expectations
and using $\Exp[X]=0$:
\[
\Exp[e^{sX}]\le\frac{b\,e^{sa}-a\,e^{sb}}{b-a}.
\]
Set $p=-a/(b-a)\in[0,1]$ and $u=s(b-a)$. Then the right side equals
$e^{-pu}(1-p+pe^u)=e^{L(u)}$ where $L(u)=-pu+\ln(1-p+pe^u)$. We have $L(0)=0$,
$L'(u)=-p+\dfrac{pe^u}{1-p+pe^u}$ so $L'(0)=0$, and
\[
L''(u)=r(1-r),\qquad r=\frac{pe^u}{1-p+pe^u}\in[0,1],
\]
so $L''(u)\le\tfrac14$ for all $u$. By Taylor's theorem with Lagrange remainder,
$L(u)=L(0)+L'(0)u+\tfrac12L''(\xi)u^2\le u^2/8=s^2(b-a)^2/8$.
\end{proof}

\begin{theorem}[Hoeffding's inequality]\label{thm:hoeffding}
Let $X_1,\dots,X_n$ be independent with $a_i\le X_i\le b_i$, and
$S_n=\sum X_i$. Then for every $t>0$,
\[
\Prob\bigl(S_n-\Exp[S_n]\ge t\bigr)\le\exp\Bigl(-\frac{2t^2}{\sum_i(b_i-a_i)^2}\Bigr),
\]
and the same bound holds for $\Prob(S_n-\Exp[S_n]\le-t)$; the two-sided probability
is at most twice the bound.
\end{theorem}
\begin{proof}
By the Chernoff method with parameter $s>0$, independence, and
Lemma~\ref{lem:hoeffding} applied to $X_i-\Exp[X_i]$ (which lies in an interval of
length $b_i-a_i$ and has mean zero),
\[
\Prob(S_n-\Exp[S_n]\ge t)\le e^{-st}\prod_i\Exp\bigl[e^{s(X_i-\Exp X_i)}\bigr]
\le\exp\Bigl(-st+\frac{s^2}8\sum_i(b_i-a_i)^2\Bigr).
\]
The exponent is minimised at $s=4t/\sum(b_i-a_i)^2$, giving
$-2t^2/\sum(b_i-a_i)^2$.
\end{proof}

\begin{example}[Polling]
Let $X_i\in\{0,1\}$ be i.i.d.\ Bernoulli$(p)$ answers and $\hat p=S_n/n$. Then
$\Prob(\abs{\hat p-p}\ge\epsilon)\le2e^{-2n\epsilon^2}$. For accuracy
$\epsilon=0.03$ with failure probability at most $5\%$ we need
$2e^{-2n(0.0009)}\le0.05$, i.e.\ $n\ge\ln(40)/(0.0018)\approx2050$ samples
(no dependence on $p$).
\end{example}

\section{Symmetry and exchangeability}

\begin{definition}[Exchangeable]
Random variables $X_1,\dots,X_n$ are \emph{exchangeable} if
$(X_{\pi(1)},\dots,X_{\pi(n)})$ has the same joint distribution as $(X_1,\dots,X_n)$
for every permutation $\pi$ of $\{1,\dots,n\}$.
\end{definition}

I.i.d.\ variables are exchangeable. So are the successive draws from an urn
\emph{without} replacement: the probability of any specific sequence of colours is
$\frac{K(K-1)\cdots}{N(N-1)\cdots}$, a product of numerators that depends only on how many
of each colour appear, not on their order.

\begin{interview}{symmetry solves many problems without computation}
If $X_1,\dots,X_n$ are exchangeable, then each $X_i$ has the same distribution, each
pair $(X_i,X_j)$ the same joint law, and so on. Example: a card dealt \emph{second}
from a shuffled deck is an ace with probability $4/52$, the same as the first.
Example (aces before kings): in a shuffled deck, what is the probability that the
first ace comes before the first king? Look only at the relative order of the $8$
cards that are aces or kings. By exchangeability the first of these $8$ is equally
likely to be any of them, so it is an ace with probability $4/8=\tfrac12$. All the
other cards are irrelevant.
\end{interview}

\begin{example}[Expected number of cards before the first ace]
The $4$ aces split the other $48$ cards into $5$ groups (before the first ace, between,
and after the last). By symmetry each of the $48$ cards is equally likely to fall in
each of the $5$ groups, so the expected size of the first group is $48/5=9.6$, and the
first ace is on average at position $9.6+1=10.6$.
\end{example}

\section*{Exercises}
\addcontentsline{toc}{section}{Exercises}

\begin{enumerate}
\item Roll a die and let $N$ be the result; then toss $N$ fair coins. Find
$\Exp[H\given N]$, $\Exp[H]$ and $\Var(H)$, where $H$ is the number of heads.
\item Show $\Var(X)=\Exp[\Var(X\given Y)]+\Var(\Exp[X\given Y])$ in the example
where $Y$ is a fair coin and $X\given Y\sim N(\pm1,1)$.
\item Find the expected number of tosses to get TTH, HTH, and HHH. Which is
largest?
\item Use Markov's inequality to show $\Prob(X\ge10)\le0.1$ if $X\ge0$ and
$\Exp[X]=1$, and show the bound can be attained.
\item Show that for a random variable with mean $0$ and variance $1$,
$\Prob(X\ge t)\le\frac1{1+t^2}$ and that this is attained for some distribution.
\item (Paley--Zygmund) Prove $\Prob(X>\theta\Exp[X])\ge(1-\theta)^2\frac{(\Exp
X)^2}{\Exp[X^2]}$ for $X\ge0$ and $0<\theta<1$. (Hint: Cauchy--Schwarz applied to
$X\indic_{\{X>\theta\Exp X\}}$.)
\item $n$ fair coins are tossed. Using Hoeffding, bound the probability that the
number of heads exceeds $n/2+\sqrt{n\ln n}$.
\item A deck is shuffled. Find the expected number of cards strictly between the
first ace and the second ace. (Symmetry.)
\item (Challenge) Prove that if $X_1,\dots,X_n$ are exchangeable with
$\Exp[X_1^2]<\infty$, then $\Cov(X_1,X_2)\ge-\Var(X_1)/(n-1)$. (Hint: compute
$\Var(\sum X_i)\ge0$.)
\end{enumerate}
